\documentclass[12pt,a4paper]{article}

% ----- Pakete -----------------------------------------------
\usepackage[T1]{fontenc}
\usepackage[utf8]{inputenc}
\usepackage[ngerman]{babel}
\usepackage{lmodern}
\usepackage{amsmath,amssymb,amsthm}
\usepackage{mathtools}
\usepackage{geometry}
\usepackage{graphicx}
\usepackage{booktabs}
\usepackage{setspace}
\usepackage{parskip}
\usepackage{hyperref}
\usepackage{titlesec}
\usepackage{fancyhdr}
\usepackage{array}
\usepackage{longtable}
\usepackage[expansion=false,protrusion=true]{microtype}
\usepackage{caption}
\usepackage{float}
\usepackage{xcolor}
\usepackage{siunitx}
\DeclareSIUnit\bar{bar}
\usepackage{bm}
\usepackage{enumitem}
\setlist[itemize,1]{label=-}

\hypersetup{colorlinks=true,linkcolor=blue!60!black,urlcolor=blue!60!black,citecolor=blue!60!black}

% ----- Seitengeometrie --------------------------------------
\geometry{margin=2.5cm}

% ----- Mathematische Umgebungen ----------------------------
\newtheorem{theorem}{Theorem}[section]
\newtheorem{lemma}[theorem]{Lemma}
\newtheorem{corollary}[theorem]{Korollar}
\newtheorem{proposition}[theorem]{Proposition}
\theoremstyle{definition}
\newtheorem{definition}[theorem]{Definition}
\newtheorem{assumption}[theorem]{Annahme}
\theoremstyle{remark}
\newtheorem{remark}[theorem]{Bemerkung}

% ----- Abkürzungen ----------------------------------------
\newcommand{\ECU}{\mathrm{ECU}}
\newcommand{\CPS}{\mathrm{CPS}}   % Central Pressure System
\newcommand{\SCS}{\mathrm{SCS}}   % Steering Control System
\newcommand{\LCS}{\mathrm{LCS}}   % Load Control System
\newcommand{\Ltire}{L_{\mathrm{tire}}}
\newcommand{\pref}{p_{\mathrm{ref}}}
\newcommand{\pact}{p_{\mathrm{act}}}
\newcommand{\Fz}{F_z}
\newcommand{\alphac}{\alpha_{\mathrm{c}}}
\newcommand{\R}{\mathbb{R}}
\newcommand{\N}{\mathbb{N}}

\title{\textbf{Maximale Lebensdauer von Gummibereifung in der Automobilindustrie}}

\author{Stephan Epp}
\date{05. März 2026}

% ===========================================================
\begin{document}

% ----- Titelseite ------------------------------------------
\maketitle

% ----- Inhaltsverzeichnis ----------------------------------
\tableofcontents
\newpage

% ===========================================================
\section{Einleitung und Motivation}

Die vorliegende Arbeit untersucht drei aufeinander abgestimmte
elektronische Steuergerätekonzepte (ECUs) zur Maximierung der
Lebensdauer von Kraftfahrzeugreifen.

\textbf{Konzept~1} beschreibt ein \emph{Zentrales Reifendruckregelungssystem} (CRDS),
das bei eingeschalteter Zündung über eine einzige Zentralbefüllungsanlage am Fahrzeug
einen lastadaptiven, temperaturkompensierten Reifeninnendruck einstellt und hält.
Es wird formal bewiesen, dass jede Abweichung vom optimalen Betriebsdruck
die Reifenlebensdauer monoton verringert.

\textbf{Konzept~2} beschreibt ein \emph{Lenksteuergerät} (SCS -- Steering Control
System), das über eine modellprädiktive Trajektorienplanung den Schräglaufwinkel
und die Querkraft auf die Reifen minimiert und damit Latsch-Mikroschlupf
sowie thermischen Verschleiß reduziert.

\textbf{Konzept~3} beschreibt ein \emph{Lastausgleichs-Steuergerät} (LCS -- Load
Control System), das die Achslastverteilung in Echtzeit überwacht, dem Fahrer
Beladungshinweise gibt und über aktive Federung die dynamische Radlast
nivelliert.

Für jedes Konzept werden physikalische Grundlagen, mathematische Modelle,
Optimierungsformulierungen und formale Beweise ausgeführt.
Abschließend wird die kombinierte Wirkung aller drei Systeme durch
ein Gesamtlebensdauermodell zusammengefasst und quantifiziert.

\textbf{Schlüsselwörter:}
Reifenlebensdauer, Reifendruck, Lenkkinematik, Lastverteilung,
Reifenverschleiß, Steuergerät, ECU, Kraftfahrzeug, Tribologie,
Regelungstechnik.

\subsection{Ausgangssituation}

Kraftfahrzeugreifen sind das einzige Bindeglied zwischen Fahrzeug und Fahrbahn.
Ihre Lebensdauer wird durch eine Vielzahl von Einflüssen bestimmt:

\begin{itemize}[nosep]
  \item Reifeninnendruck (Unter- bzw.\ Überdruck führen zu ungleichmäßigem Abrieb),
  \item Lenkgeometrie und Schräglaufwinkel (erzeugen Seitenkräfte und Latschmikroschlupf),
  \item Radlasten und deren Verteilung auf alle Achsen (Ãœberlastung einzelner Reifen),
  \item Fahrgeschwindigkeit und Fahrstil,
  \item Umgebungstemperatur und Fahrbahnoberflächenbeschaffenheit.
\end{itemize}

Die wirtschaftliche Bedeutung ist erheblich: In der EU werden jährlich ca.\ 450
Millionen PKW-Reifen verbraucht~\cite{ETRMA2023}. Ein Anstieg der durchschnittlichen
Reifenlebensdauer um nur 15\,\% würde den jährlichen Ressourcenverbrauch um ca.\ 67,5
Millionen Reifen verringern. Hinzu kommen Effizienzgewinne durch geringeren Rollwiderstand
bei optimalem Reifendruck~\cite{Sandberg2011}.

\subsection{Zielsetzung der Arbeit}

Diese Arbeit verfolgt drei komplementäre Ziele:

\begin{enumerate}
  \item \textbf{Reifendruckregelung:} Formal-mathematischen Nachweis erbringen, dass
        eine zentrale, steuergerätgeführte Druckregelung bei eingeschalteter Zündung
        die Reifenlebensdauer gegenüber manueller Druckpflege maximiert.
  \item \textbf{Lenksteuerung:} Nachweis, dass ein modellprädiktiver Lenkungsregler
        die reifenschädigenden Seitenkräfte und Schräglaufwinkel über den gesamten
        Fahrzyklus minimiert.
  \item \textbf{Lastausgleich:} Nachweis, dass eine gleichmäßige dynamische Radlastverteilung
        durch ein aktives Fahrwerk- und Beladungshinweissystem die Gesamtreifenlebensdauer
        eines Reifensatzes maximiert.
\end{enumerate}

\subsection{Aufbau der Arbeit}

Abschnitt~\ref{sec:grundlagen} legt die physikalischen und tribologischen Grundlagen
des Reifenverschleißes dar. Die Abschnitte~\ref{sec:crds}, \ref{sec:scs} und
\ref{sec:lcs} behandeln die drei ECU-Konzepte mit vollständigem mathematischen Apparat
und formalen Beweisen. Abschnitt~\ref{sec:gesamtmodell} fasst das Gesamtmodell zusammen.
Abschnitt~\ref{sec:schluss} gibt Schlussfolgerungen und einen Ausblick.
Das Quellenverzeichnis am Ende der Arbeit ist vollständig und inline eingebettet.

% ===========================================================
\section{Physikalische und Tribologische Grundlagen}
\label{sec:grundlagen}

\subsection{Reifen-Fahrbahn-Kontaktmodell}

\subsubsection{Latschgeometrie und Kontaktdruck}

Der Reifenlatsch (Kontaktfläche Reifen–Fahrbahn) kann für einen pneumatischen
Reifen unter der vereinfachenden Annahme einer rechteckigen Fläche modelliert
werden~\cite{Pacejka2012}:

\begin{equation}
  A_{\mathrm{latsch}} = \frac{\Fz}{p_i}
  \label{eq:latsch}
\end{equation}

wobei $\Fz$ die Radlast [\si{\newton}] und $p_i$ der Reifeninnendruck
[\si{\pascal}] sind. Der mittlere Kontaktdruck $\sigma_0$ beträgt dann:

\begin{equation}
  \sigma_0 = \frac{\Fz}{A_{\mathrm{latsch}}} = p_i
  \label{eq:kontaktdruck}
\end{equation}

Diese Näherung ist für Standardreifen bis ca.\ $\pm 20$\,\% des Nenndrucks
valide~\cite{Gent2006}. Eine präzisere Modellierung berücksichtigt die
Karkassensteifigkeit $k_c$:

\begin{equation}
  A_{\mathrm{latsch,korr}} = \frac{\Fz}{p_i + k_c / r_d}
  \label{eq:latsch_korr}
\end{equation}

wobei $r_d$ den dynamischen Reifenrollradius bezeichnet.

\subsubsection{Schräglauftheorie nach Pacejka}

Der laterale Kraftaufbau im Reifen folgt dem \emph{Magic Formula}-Modell
von Pacejka, siehe dazu \cite{Pacejka2012}. Für die Seitenkraft $F_y$ als Funktion des
Schräglaufwinkels $\alpha$ gilt:

\begin{equation}
  F_y(\alpha) = D \cdot \sin\!\bigl[C \cdot \arctan\!\bigl(B\alpha - E(B\alpha
    - \arctan(B\alpha))\bigr)\bigr]
  \label{eq:magic_formula}
\end{equation}

Hierbei sind:
\begin{itemize}[nosep]
  \item $B$ -- Steifigkeitsfaktor,
  \item $C$ -- Formfaktor (typisch $C \approx 1{,}3$ für Seitenkräfte),
  \item $D$ -- Spitzenwert der Seitenkraft,
  \item $E$ -- Krümmungsfaktor.
\end{itemize}

\subsection{Reifenverschleißmodell}

\subsubsection{Archard-Modell für Reifenabrieb}

Reifenabrieb ist primär Adhäsions- und Abrasivverschleiß. Das klassische
Archard-Verschleiß\-gesetz~\cite{Archard1953} lautet:

\begin{equation}
  \dot{V}_w = k_w \cdot \frac{\Fz \cdot v_s}{H}
  \label{eq:archard}
\end{equation}

wobei $\dot{V}_w$ das Volumenverschleißrate [\si{\metre\cubed\per\second}],
$k_w$ der Verschleißkoeffizient (Archard-Konstante), $v_s$ die
Schlupfgeschwindigkeit [\si{\metre\per\second}] und $H$ die
Reifengummihärte [\si{\pascal}] sind.

\subsubsection{Energetisches Verschleißmodell}

Moderner und für Reifen präziser ist das energetische Modell nach
Schallamach und Grosch \cite{Schallamach1971}, \cite{Grosch1963}:

\begin{equation}
  \dot{m}_w = k_E \cdot p \cdot v_s
  \label{eq:energetisch}
\end{equation}

wobei $\dot{m}_w$ die Masseverschleißrate, $k_E$ der energetische
Verschleißkoeffizient, $p$ der lokale Kontaktdruck und $v_s$ die lokale
Schlupfgeschwindigkeit sind.

\subsubsection{Reifenlebensdauer als integrierte Größe}

Die Gesamtlebensdauer eines Reifens $\Ltire$ wird definiert als der
Zeitpunkt, an dem die Profiltiefe $h(t)$ den gesetzlich vorgeschriebenen
Mindest\-wert $h_{\min}$ (in Deutschland \SI{1,6}{\milli\metre}) unterschreitet:

\begin{equation}
  \Ltire = \min\!\left\{t > 0 : h(t) \leq h_{\min}\right\}
  \label{eq:life_def}
\end{equation}

Die Profiltiefe entwickelt sich als:

\begin{equation}
  h(t) = h_0 - \int_0^t \dot{h}(\tau)\,\mathrm{d}\tau
  \label{eq:profiltiefe}
\end{equation}

wobei $h_0$ die Anfangsprofiltiefe ist und $\dot{h}(\tau)$ die momentane
Abtragrate der Profiltiefe bezeichnet. Letztere ist proportional zur
Schlupfleistungsdichte:

\begin{equation}
  \dot{h}(t) = c_h \cdot \int_{A_{\mathrm{latsch}}} p(\xi)\,v_s(\xi,t)\,\mathrm{d}\xi
  \label{eq:hpunkt}
\end{equation}

mit der reifenspezifischen Konstante $c_h > 0$.

\subsection{Temperaturabhängigkeit}

Die Reifentemperatur $T_r$ beeinflusst $k_E$ exponentiell (WLF-Gesetz~\cite{Williams1955}):

\begin{equation}
  k_E(T_r) = k_{E,\mathrm{ref}}\cdot\exp\!\left(\frac{C_1(T_r - T_{\mathrm{ref}})}
    {C_2 + T_r - T_{\mathrm{ref}}}\right)
  \label{eq:wlf}
\end{equation}

mit den WLF-Konstanten $C_1 \approx 17{,}44$, $C_2 \approx 51{,}6$\,K
(bezogen auf $T_{\mathrm{ref}} = T_g + 50$\,K, $T_g$ = Glasübergangstemperatur).

% ===========================================================
\section{Konzept 1: Zentrales Reifendruckregelungssystem (CRDS)}
\label{sec:crds}

\subsection{Systemarchitektur}

Das CRDS umfasst:
\begin{itemize}
  \item Eine zentrale Hochdruckpumpe und einen Luftbehälter ($p_{\max}$
        $\approx$ \SI{10}{\bar}) mit elektromagnetischen Sperrventilen,
  \item Je einen Drehübertrager (rotierender Schlauchanschluss) pro Achse,
  \item Raddrucksensoren an allen vier Rädern,
  \item Temperatursensoren im Reifeninnern oder Näherungsschätzung via
        Wärmemodell,
  \item Einen zentralen Mikrokontroller (CRDS-ECU) mit CAN-Bus-Anbindung,
  \item Ablassventile zum kontrollierten Druckablassen bei Ãœberdruck.
\end{itemize}

Das System ist nur aktiv, wenn Zündung EIN und Fahrzeuggeschwindigkeit
$v < v_{\mathrm{threshold}}$ oder der Druck außerhalb des Toleranzbandes liegt.

\subsection{Optimaler Reifeninnendruck}

\subsubsection{Druckabhängige Verschleißrate}

Aus dem Latschmodell (Gl.~\ref{eq:latsch}) und dem Verschleißmodell
(Gl.~\ref{eq:energetisch}) folgt die druckabhängige Verschleißrate.
Für einen symmetrischen Latsch (keine Schräglaufkraft) gilt mit dem
Rollschlupf $\kappa$ (typisch $|\kappa| < 1$\,\%):

\begin{equation}
  \dot{h}(p_i) = c_h \cdot \Fz \cdot |\kappa| \cdot f_{\mathrm{dist}}(p_i)
  \label{eq:hdot_pi}
\end{equation}

wobei $f_{\mathrm{dist}}(p_i)$ die druckabhängige Verteilungsfunktion der
Kontaktspannung beschreibt. Bei Unterdruck bildet sich eine Randlasthöhung
(``Trogform''), bei Überdruck eine Mittenlasthöhung (``Kugelform'')~\cite{Gent2006}:

\begin{equation}
  f_{\mathrm{dist}}(p_i) = 1 + \beta \left(\frac{p_i - \pref}{\pref}\right)^2,
  \quad \beta > 0
  \label{eq:fdist}
\end{equation}

Das Minimum liegt bei $p_i = \pref$ mit $f_{\mathrm{dist}}(\pref) = 1$.

\begin{theorem}[Optimaler Reifendruck minimiert Verschleiß]
\label{thm:pressure}
Sei $\dot{h}(p_i)$ gemäß Gl.~\eqref{eq:hdot_pi}--\eqref{eq:fdist} mit
$\beta > 0$. Dann gilt:
\begin{equation}
  \dot{h}(\pref) \leq \dot{h}(p_i) \quad \forall\, p_i > 0
  \label{eq:thm_pressure}
\end{equation}
und $\Ltire(p_i)$ ist maximal bei $p_i = \pref$.
\end{theorem}

\begin{proof}
Da $c_h > 0$, $\Fz > 0$ und $|\kappa| \geq 0$ sind konstante positive Faktoren,
genügt es, $f_{\mathrm{dist}}(p_i)$ zu minimieren.
Es gilt:
\begin{equation}
  f_{\mathrm{dist}}(p_i) = 1 + \beta \underbrace{\left(\frac{p_i - \pref}{\pref}\right)^2}_{\geq\, 0}
  \geq 1 = f_{\mathrm{dist}}(\pref)
\end{equation}
da $\beta > 0$ und das Quadrat nicht-negativ ist. Das Minimum wird genau bei
$p_i = \pref$ angenommen. Aus \eqref{eq:profiltiefe} folgt:
\begin{equation}
  h(t; p_i) = h_0 - \int_0^t \dot{h}(p_i)\,\mathrm{d}\tau
            = h_0 - \dot{h}(p_i) \cdot t
  \quad\text{(für konstantes $p_i$)}
\end{equation}
Aus $h(t; p_i) = h_{\min}$ folgt:
\begin{equation}
  \Ltire(p_i) = \frac{h_0 - h_{\min}}{\dot{h}(p_i)}
\end{equation}
Da $\dot{h}(p_i) \geq \dot{h}(\pref)$, gilt:
\begin{equation}
  \Ltire(p_i) = \frac{h_0 - h_{\min}}{\dot{h}(p_i)}
              \leq \frac{h_0 - h_{\min}}{\dot{h}(\pref)} = \Ltire(\pref)
\end{equation}
Also ist $\Ltire$ maximal bei $p_i = \pref$.
\end{proof}

\subsection{Lastabhängige Druckkorrektur}

Der optimale Druck ist last- und temperaturabhängig. Die Referenzgröße lautet:

\begin{equation}
  \pref(\Fz, T_r) = p_0 \cdot \frac{\Fz}{F_{z,\mathrm{ref}}} \cdot
    \frac{T_r + 273{,}15}{T_{r,\mathrm{ref}} + 273{,}15}
  \label{eq:pref_adapt}
\end{equation}

Dabei ist $p_0$ der Referenzdruck bei Nennlast $F_{z,\mathrm{ref}}$ und
Referenztemperatur $T_{r,\mathrm{ref}}$ (üblicherweise \SI{20}{\celsius}).
Die Temperaturkorrektur folgt dem idealen Gasgesetz (Gay-Lussac).

\subsection{Regelungsgesetz des CRDS-ECU}

Der diskrete Regelkreis (Abtastintervall $T_s$) lautet:

\begin{align}
  e_k &= \pref\!\left(\Fz^{(k)}, T_r^{(k)}\right) - \pact^{(k)} \label{eq:crds_e}\\
  u_k &= K_P \, e_k + K_I \sum_{j=0}^{k} e_j \, T_s + K_D \frac{e_k - e_{k-1}}{T_s}
       \label{eq:crds_pid}
\end{align}

mit der Stellgröße $u_k \in \{-1, 0, +1\}$ (Ablassen, Halten, Befüllen).
Der Beweis der Stabilität und Stationärgenauigkeit dieses PID-Reglers
folgt aus dem Nyquist-Kriterium für das diskrete System~\cite{Astrom1997}.

\subsubsection{Stabilitätsbeweis des Druckreglers}

\begin{theorem}[Stabilität des CRDS-Reglers]
\label{thm:crds_stability}
Der PI-Druckregler (Gl.~\eqref{eq:crds_e}--\eqref{eq:crds_pid}) konvergiert
für das PT1-Streckenmodell der pneumatischen Strecke gegen den stationären
Sollwert $\pref$, wenn die Reglerverstärkungen $K_P$, $K_I$ so gewählt sind,
dass alle Pole des geschlossenen Kreises innerhalb des Einheitskreises liegen.
\end{theorem}

\begin{proof}
Das Streckenmodell (Befüllung durch ein Magnetventil in ein
Volumen $V$ mit Widerstand $R_v$) ist ein PT1-Glied:
\begin{equation}
  G(s) = \frac{K_{str}}{1 + T_{str} s}, \quad
  T_{str} = R_v \cdot V, \quad K_{str} > 0
\end{equation}
Der PI-Regler hat die Ãœbertragungsfunktion:
\begin{equation}
  C(s) = K_P + \frac{K_I}{s} = \frac{K_P s + K_I}{s}
\end{equation}
Das offene System:
\begin{equation}
  L(s) = C(s) G(s) = \frac{K_{str}(K_P s + K_I)}{s(1 + T_{str} s)}
\end{equation}
Das charakteristische Polynom des geschlossenen Kreises:
\begin{equation}
  1 + L(s) = 0 \;\Rightarrow\; T_{str} s^2 + (1 + K_{str} K_P) s + K_{str} K_I = 0
\end{equation}
Nach dem Hurwitz-Kriterium sind alle Koeffizienten positiv, wenn
$K_P > -1/K_{str}$ (in der Praxis trivial erfüllt für $K_P > 0$) und $K_I > 0$.
Damit sind beide Pole in der linken Halbebene. Für das diskrete System wird
mit der Tustin-Transformation $s = \frac{2}{T_s}\frac{z-1}{z+1}$ die entsprechende
Stabilitätsbedingung im Einheitskreis sichergestellt.
\end{proof}

\subsection{Energiebetrachtung}

Die Energiekosten der Druckregelung sind gering. Der Arbeitsbedarf zur
Befüllung eines Reifens um $\Delta p$ bei Volumen $V_r$ beträgt näherungsweise:

\begin{equation}
  W_{\text{Befüllung}} = V_r \cdot \Delta p \cdot \ln\!\frac{p_{\mathrm{final}}}{p_{\mathrm{init}}}
  \approx V_r \cdot \Delta p
  \quad\text{für kleine $\Delta p$}
\end{equation}

Für $V_r \approx$ \SI{20}{\litre}, $\Delta p \approx$ \SI{0{,}1}{\bar} ergibt sich
$W \approx$ \SI{200}{\joule}, was bei einer Kompressorleistung von \SI{200}{\watt}
einer Betriebsdauer von \SI{1}{\second} entspricht -- energetisch vernachlässigbar.

\subsection{Sicherheitskonzept}

Das CRDS muss folgende Sicherheitsfunktionen implementieren:

\begin{enumerate}
  \item \textbf{Redundante Drucksensoren:} Mindestens zwei unabhängige Drucksensoren
        pro Rad mit Plausibilitätsprüfung.
  \item \textbf{Fail-Safe:} Bei Sensorausfall wird der letzte valide Druckwert
        gehalten; der Fahrer erhält eine Warnung.
  \item \textbf{Ãœberdruck-Schutzventil:} Mechanisches Ãœberdruckventil bei
        $p > p_{\max,\mathrm{safe}}$.
  \item \textbf{Aktivierung nur bei Zündung EIN:} Verhindert Entladung der
        12-V-Batterie bei stehendem Fahrzeug.
\end{enumerate}

% ===========================================================
\section{Konzept 2: Lenksteuergerät mit Reifenschonender Trajektorienplanung (SCS)}
\label{sec:scs}

\subsection{Motivation: Seitenkraft als Hauptverschleißursache}

Seitenkräfte (Querkräfte) am Reifen erzeugen Latschmikroschlupf und sind
nach Rollwiderstand und Längsschlupf die dominierende Verschleißursache~\cite{Niskanen2014}.
Die seitliche Verschleißrate ist:

\begin{equation}
  \dot{h}_{\mathrm{lat}} = c_y \cdot |F_y(\alpha)| \cdot |\alpha|
  \label{eq:wear_lat}
\end{equation}

Da $F_y(\alpha)$ für kleine $\alpha$ näherungsweise linear ist ($F_y \approx C_\alpha \alpha$,
mit der Schräglaufsteifigkeit $C_\alpha$):

\begin{equation}
  \dot{h}_{\mathrm{lat}} \approx c_y \cdot C_\alpha \cdot \alpha^2
  \label{eq:wear_lat_lin}
\end{equation}

\subsection{Kinematisches Fahrzeugmodell}

\subsubsection{Einspurmodell (Bicycle Model)}

Das vereinfachte Einspurmodell~\cite{Rajamani2006} beschreibt die Fahrdynamik
durch den Schwimmwinkel $\beta$, die Gierrate $\dot{\psi}$ und den Lenkwinkel $\delta$:

\begin{align}
  m v (\dot{\beta} + \dot{\psi}) &= F_{y,v} + F_{y,h} \label{eq:bicycle1}\\
  I_z \ddot{\psi} &= l_v F_{y,v} - l_h F_{y,h} \label{eq:bicycle2}
\end{align}

wobei:
\begin{itemize}[nosep]
  \item $m$ -- Fahrzeugmasse,
  \item $v$ -- Fahrzeuggeschwindigkeit,
  \item $l_v, l_h$ -- Abstände Schwerpunkt--Vorder-/Hinterachse,
  \item $F_{y,v}, F_{y,h}$ -- Seitenkräfte an Vorder-/Hinterachse,
  \item $I_z$ -- Massenträgheitsmoment um Hochachse.
\end{itemize}

Die Schräglaufwinkel der Achsen betragen:

\begin{align}
  \alpha_v &= \delta - \beta - \frac{l_v \dot{\psi}}{v} \label{eq:alpha_v}\\
  \alpha_h &= -\beta + \frac{l_h \dot{\psi}}{v} \label{eq:alpha_h}
\end{align}

\subsection{Optimaler Lenkwinkel zur Minimierung des Reifenverschleißes}

\subsubsection{Optimierungsproblem}

Für eine gegebene Trajektorie (Zielkrümmung $\kappa_{\mathrm{ref}}$, Geschwindigkeit $v$)
wird der Lenkwinkel $\delta$ so gewählt, dass der Schräglaufwinkel minimiert wird.

Die Gesamtverschleißrate aller vier Reifen lautet:

\begin{equation}
  J_w(\delta) = \sum_{i \in \{vl, vr, hl, hr\}} \dot{h}_i(\alpha_i(\delta))
  = c_y C_\alpha \sum_{i} \alpha_i^2(\delta)
  \label{eq:wear_total}
\end{equation}

Das SCS-Optimierungsproblem:

\begin{equation}
  \delta^*_k = \arg\min_{\delta \in [\delta_{\min}, \delta_{\max}]}
    J_w(\delta) \quad \text{subject to: } \kappa(x_k, \delta_k) = \kappa_{\mathrm{ref},k}
  \label{eq:scs_opt}
\end{equation}

\begin{theorem}[Existenz und Eindeutigkeit des optimalen Lenkwinkels]
\label{thm:scs_unique}
Für gegebene $\kappa_{\mathrm{ref}}$, $v > 0$ und ein lineares Reifenmodell besitzt
das Optimierungsproblem \eqref{eq:scs_opt} eine eindeutige Lösung $\delta^*$.
\end{theorem}

\begin{proof}
Im linearen Bereich sind die Schräglaufwinkel $\alpha_v, \alpha_h$ affin-linear
in $\delta$ (vgl.\ \eqref{eq:alpha_v}, \eqref{eq:alpha_h}). Das Quadrat ist konvex;
eine endliche Summe konvexer Funktionen ist konvex. Die Nebenbedingung
$\kappa(x_k, \delta_k) = \kappa_{\mathrm{ref}}$ ist für das Einspurmodell eine
Gleichungsnebenbedingung und linear in $\delta$ für kleine Winkel. Damit ist
das Problem ein konvex beschränktes QP (quadratisches Programm) in einer
Variablen -- die Lösung existiert und ist eindeutig.
\end{proof}

\subsubsection{Analytische Lösung für symmetrisches Fahrzeug}

Für den Sonderfall $l_v = l_h = l/2$ (symmetrischer Radstand) und
$C_{\alpha,v} = C_{\alpha,h} = C_\alpha$ ergibt sich die Ackermann-Bedingung
als optimaler Lenkwinkel:

\begin{equation}
  \delta^*_{\mathrm{Ackermann}} = \arctan\!\frac{l}{\rho}
  \approx \frac{l}{\rho}
  \label{eq:ackermann}
\end{equation}

wobei $\rho$ der Kurvenradius ist. Dies minimiert beide Schräglaufwinkel
simultaneously auf Null, was dem globalen Minimum von $J_w$ entspricht.

\begin{corollary}[Ackermann-Lenkung als Verschleißminimum]
Die ideale Ackermann-Lenkgeometrie minimiert die reifenverschleißende
Seitenführungskraft für stationäres Kurvenfahren.
\end{corollary}

\subsection{Modellprädiktive Lenkungsregelung (MPC)}

\subsubsection{Systemzustand und Prädiktionshorizont}

Der Zustandsvektor sei $\mathbf{x} = [\beta, \dot{\psi}, v]^T$, die Stellgröße
$u = \delta$ (Lenkwinkel). Das diskrete Zustandsraummodell lautet:

\begin{equation}
  \mathbf{x}_{k+1} = \mathbf{A}\,\mathbf{x}_k + \mathbf{B}\,u_k
  \label{eq:state_space}
\end{equation}

Der MPC-Regler löst über den Prädiktionshorizont $N$ das Optimierungsproblem:

\begin{equation}
  \min_{\{u_k, \ldots, u_{k+N-1}\}} \sum_{j=k}^{k+N-1}
    \left[(\mathbf{x}_j - \mathbf{x}_{\mathrm{ref},j})^T \mathbf{Q}
    (\mathbf{x}_j - \mathbf{x}_{\mathrm{ref},j}) + u_j^T R\, u_j\right]
  \label{eq:mpc}
\end{equation}

mit den Gewichtsmatrizen $\mathbf{Q} \succeq 0$ (Zustandskosten) und
$R > 0$ (Stellgrößenkosten, penalisiert zu starkes Lenken).

\subsubsection{Reifenverschleiß als Optimierungsziel}

Durch entsprechende Wahl von $\mathbf{Q}$ -- mit hohem Gewicht auf den
Schräglaufwinkeln $\alpha_v, \alpha_h$ -- wird der Verschleißterm in das
MPC-Kostenfunktional integriert:

\begin{equation}
  Q_{\alpha} = \text{diag}(q_\beta, q_{\dot\psi}, q_v), \quad
  q_\beta \gg q_{\dot\psi},\; q_v
  \label{eq:q_weights}
\end{equation}

\begin{theorem}[MPC-Optimierung reduziert Reifenverschleiß]
\label{thm:mpc_wear}
Sei der Reifenverschleiß eine stetig differenzierbare, konvexe Funktion der
Schräglaufwinkel. Dann konvergiert die MPC-Lösung für $N \to \infty$ gegen
die Trajektorie mit minimalem Gesamtverschleiß unter Einhaltung der
Trajektorienvorgaben.
\end{theorem}

\begin{proof}
Für konvexe Kostenfunktionale und lineare Systemdynamik ist das MPC-Problem
ein konvexes QP. Für $N \to \infty$ nähert sich die MPC-Lösung der
\emph{infinite-horizon} optimalen Regelung an (LQR-Äquivalenz)~\cite{Mayne2000}.
Das Gesamtlebensdauerintegral $\int_0^T J_w(t)\,\mathrm{d}t$ wird minimiert,
da jedes endliche Zeitsegment durch die MPC-Lösung sub-optimal (für endliches $N$)
bzw.\ optimal (für $N \to \infty$) behandelt wird. Damit ist die resultierende
Trajektorie verschleißoptimal.
\end{proof}

\subsection{Reifenverschleiß-Ausgleich durch Spurwinkeldämpfung}

In Geradeausfahrt kann das SCS durch kleine Lenkkorrekturen die Abnutzung
über die Reifenbreite verteilen (``Wear Leveling''). Das SCS-ECU überwacht
die kumulierte Schlupfenergie $E_s = \int_0^t F_y \cdot v_s \,\mathrm{d}\tau$
pro Reifensegment und gibt Lenkkorrekturen, die die Kontaktpunktverteilung
optimieren.

% ===========================================================
\section{Konzept 3: Lastausgleichs-Steuergerät mit Fahrerhinweissystem (LCS)}
\label{sec:lcs}

\subsection{Einfluss der Radlastverteilung auf die Reifenlebensdauer}

\subsubsection{Nichtlinearer Zusammenhang zwischen Last und Verschleiß}

Das Verschleißgesetz (Gl.~\ref{eq:energetisch}) ist \emph{nichtlinear} in der
Radlast $\Fz$, wenn man die Schlupfgeschwindigkeit berücksichtigt:

\begin{equation}
  \dot{h}(\Fz) = c_h \cdot \Fz^n, \quad n > 1
  \label{eq:wear_fz}
\end{equation}

Empirisch gilt für Autobahnreifen $n \approx 1{,}5$ bis $2{,}0$~\cite{Niskanen2014,
Williams1955}.

\begin{theorem}[Gleichmäßige Lastverteilung maximiert Reifensatzlebensdauer]
\label{thm:load_even}
Sei die Verschleißrate jedes Reifens $i$ gemäß $\dot{h}_i = c_h F_{z,i}^n$,
$n > 1$, und die Gesamtlast $\sum_{i=1}^4 F_{z,i} = F_{\mathrm{ges}} = \mathrm{const}$.
Die Lebensdauer des Reifensatzes sei definiert als:
\begin{equation}
  \Ltire^{\mathrm{set}} = \min_i L_{\mathrm{tire},i}
                        = \min_i \frac{h_0 - h_{\min}}{c_h F_{z,i}^n}
  \label{eq:life_set}
\end{equation}
Dann wird $\Ltire^{\mathrm{set}}$ maximiert bei gleichmäßiger Lastverteilung
$F_{z,i} = F_{\mathrm{ges}} / 4$ für alle $i$.
\end{theorem}

\begin{proof}
Da $L_{\mathrm{tire},i} = K / F_{z,i}^n$ mit $K = (h_0 - h_{\min}) / c_h > 0$, ist
$\Ltire^{\mathrm{set}} = K / \max_i(F_{z,i}^n) = K / (\max_i F_{z,i})^n$.
Das Minimum der Lebensdauern wird also durch den höchstbelasteten Reifen
bestimmt. Da $f(x) = x^n$ für $n > 1$ strikt konvex und monoton steigend ist,
gilt:
\begin{equation}
  \max_i F_{z,i} \geq \frac{F_{\mathrm{ges}}}{4}
\end{equation}
mit Gleichheit genau dann, wenn $F_{z,1} = F_{z,2} = F_{z,3} = F_{z,4} = F_{\mathrm{ges}}/4$.
Damit ist $\Ltire^{\mathrm{set}} = K \cdot (4/F_{\mathrm{ges}})^n$ bei
gleichmäßiger Verteilung maximal.
\end{proof}

\subsubsection{Konvexitätsargument für allgemeine Verschleißfunktionen}

Das Ergebnis lässt sich verallgemeinern:

\begin{proposition}[Jensensche Ungleichung für Reifenverschleiß]
\label{prop:jensen}
Sei $\phi : \R_{>0} \to \R_{>0}$ strikt konvex und $\phi' > 0$.
Für den mittleren Verschleiß gilt:
\begin{equation}
  \frac{1}{4}\sum_{i=1}^4 \phi(F_{z,i}) \geq \phi\!\left(\frac{1}{4}\sum_{i=1}^4 F_{z,i}\right)
    = \phi\!\left(\frac{F_{\mathrm{ges}}}{4}\right)
  \label{eq:jensen}
\end{equation}
mit Gleichheit genau bei $F_{z,1} = \cdots = F_{z,4}$.
\end{proposition}

\begin{proof}
Direkte Anwendung der Jensenschen Ungleichung~\cite{Jensen1906} für strikt
konvexe Funktionen.
\end{proof}

\subsection{Radlastmessung und Schätzung}

Das LCS-ECU ermittelt die Radlasten aus:
\begin{itemize}
  \item \textbf{Federwegssensoren:} $F_{z,i} = k_{f,i} \cdot x_{f,i} + d_{f,i} \cdot \dot{x}_{f,i}$
        (linear-viskoelastisches Federmodell),
  \item \textbf{Beschleunigungssensoren} (zur Filterung dynamischer Anteile),
  \item \textbf{Gewichtssensoren} in den Sitzschienen (Insassengewicht),
  \item \textbf{Laderaum-Drucksensoren} oder kamerabasierte Volumen/Schwerpunktschätzung.
\end{itemize}

\subsubsection{Kalman-Filter für Radlastschätzung}

Zur Rauschunterdrückung verwendet das LCS-ECU einen Kalman-Filter~\cite{Kalman1960}
mit dem Zustandsvektor:

\begin{equation}
  \mathbf{z}_k = \left[F_{z,1}, F_{z,2}, F_{z,3}, F_{z,4}, F_{\mathrm{Zuladung}}\right]^T
\end{equation}

Der Prädiktionsschritt:
\begin{align}
  \hat{\mathbf{z}}_{k|k-1} &= \mathbf{A}_z \hat{\mathbf{z}}_{k-1|k-1}\\
  \mathbf{P}_{k|k-1} &= \mathbf{A}_z \mathbf{P}_{k-1|k-1} \mathbf{A}_z^T + \mathbf{Q}_z
\end{align}

Der Korrekturschritt:
\begin{align}
  \mathbf{K}_k &= \mathbf{P}_{k|k-1}\mathbf{H}^T
    (\mathbf{H}\mathbf{P}_{k|k-1}\mathbf{H}^T + \mathbf{R}_z)^{-1}\\
  \hat{\mathbf{z}}_{k|k} &= \hat{\mathbf{z}}_{k|k-1} + \mathbf{K}_k
    (\mathbf{y}_k - \mathbf{H}\hat{\mathbf{z}}_{k|k-1})\\
  \mathbf{P}_{k|k} &= (\mathbf{I} - \mathbf{K}_k \mathbf{H})\mathbf{P}_{k|k-1}
\end{align}

\subsection{Fahrerhinweissystem beim Beladen}

\subsubsection{Lastverteilungs-Index}

Das LCS-ECU berechnet den \emph{Lastverteilungs-Ungleichgewichts-Index} (LUI):

\begin{equation}
  \mathrm{LUI} = \frac{\max_i F_{z,i} - \min_i F_{z,i}}{F_{\mathrm{ges}} / 4}
  \label{eq:lui}
\end{equation}

$\mathrm{LUI} = 0$ entspricht perfekter Gleichverteilung; Grenzwerte:
\begin{itemize}[nosep]
  \item $\mathrm{LUI} < 0{,}1$: Grün (optimal),
  \item $0{,}1 \leq \mathrm{LUI} < 0{,}3$: Gelb (Hinweis empfohlen),
  \item $\mathrm{LUI} \geq 0{,}3$: Rot (Warnung, Reifenlebensdauer stark beeinträchtigt).
\end{itemize}

\subsubsection{Beladungshinweis-Algorithmus}

Beim Einschalten der Zündung führt das LCS-ECU folgende Routine aus:

\begin{enumerate}
  \item Messung aller vier Radlasten $F_{z,i}^{(0)}$.
  \item Berechnung von LUI und dem optimalen Beladungsschwerpunkt
        $\mathbf{r}^* = [x^*, y^*]^T$:
        \begin{equation}
          x^* = \frac{l_v}{2}, \quad
          y^* = 0
          \quad \text{(Fahrzeugmitte)}
          \label{eq:optimal_load_center}
        \end{equation}
  \item Vergleich des aktuellen Schwerpunkts der Zuladung:
        \begin{equation}
          x_{\mathrm{act}} = \frac{l_h \sum_v F_{z,i} - l_v \sum_h F_{z,i}}{F_{\mathrm{Zuladung}}}
          \label{eq:actual_cog}
        \end{equation}
  \item Ausgabe eines sprachlichen Hinweises an den Fahrer über das Display/Sprachsystem:
        z.\,B.\ ``Gepäck bitte \SI{20}{\centi\metre} weiter nach vorne und
        \SI{10}{\centi\metre} zur Mitte verschieben.''
\end{enumerate}

\subsection{Aktives Fahrwerk als ergänzende Maßnahme}

Neben der Beladungsoptimierung kann ein aktives (semi-aktives) Fahrwerksystem die
dynamischen Radlastschwankungen $\Delta F_{z,i}(t)$ reduzieren:

\begin{equation}
  F_{z,i}(t) = F_{z,i}^{(0)} + \Delta F_{z,i}(t), \quad
  \Delta F_{z,i}(t) = -c_a(t)\,z_i(t) - d_a(t)\,\dot{z}_i(t)
\end{equation}

wobei $c_a, d_a$ die aktiv geregelten Feder- und Dämpferkoeffizienten sind.

\begin{theorem}[Dynamische Radlastminimierung durch aktives Fahrwerk]
\label{thm:active_suspension}
Sei die mittlere quadratische Radlastschwankung:
\begin{equation}
  J_\Delta = \sum_{i=1}^4 \mathbb{E}\!\left[(\Delta F_{z,i})^2\right]
\end{equation}
Ein optimaler $\mathcal{H}_2$-Regler für das Fahrwerkssystem minimiert $J_\Delta$
und reduziert damit die nichtlineare Verschleißverstärkung aus Gl.~\eqref{eq:wear_fz}.
\end{theorem}

\begin{proof}
Für Stochastik mit weißem Rauschen als Straßenanregung ist die $\mathcal{H}_2$-Optimierung
das LQG-Problem (Linear Quadratic Gaussian)~\cite{Anderson1990}. Die Existenz eines
optimalen Reglers ist garantiert, wenn das System steuerbar und beobachtbar ist
(erfüllt für Standardfahrwerksmodelle). Die Minimierung von $J_\Delta$ verringert
$\mathbb{E}[F_{z,i}^n] = (F_{z,i}^{(0)})^n + \binom{n}{2}(F_{z,i}^{(0)})^{n-2}\mathbb{E}[(\Delta F_{z,i})^2]
+ \mathcal{O}(\mathbb{E}[(\Delta F_{z,i})^3])$ für $n \geq 2$, womit die mittlere
Verschleißrate sinkt.
\end{proof}

% ===========================================================
\section{Gesamtmodell: Kombinierte Lebensdauersteigerung}
\label{sec:gesamtmodell}

\subsection{Multiplikatives Verschleißmodell}

Die Einzelbeiträge der drei Konzepte können als multiplikative Faktoren
auf die Basisverschleißrate modelliert werden:

\begin{equation}
  \dot{h}_{\mathrm{ges}}(t) = \dot{h}_{\mathrm{basis}} \cdot
    \underbrace{\Phi_P(p_i(t))}_{\text{CRDS}} \cdot
    \underbrace{\Phi_S(\alpha(t))}_{\text{SCS}} \cdot
    \underbrace{\Phi_L(\Fz(t))}_{\text{LCS}}
  \label{eq:total_wear}
\end{equation}

mit den dimensionslosen Verschleißfaktoren:

\begin{align}
  \Phi_P(p_i) &= f_{\mathrm{dist}}(p_i) = 1 + \beta\left(\frac{p_i - \pref}{\pref}\right)^2
                \geq 1 \label{eq:phi_p}\\
  \Phi_S(\alpha) &= 1 + \gamma \alpha^2 \geq 1 \label{eq:phi_s}\\
  \Phi_L(\Fz) &= \left(\frac{\Fz}{F_{\mathrm{ref}}}\right)^{n-1} \geq 1
    \quad\text{für } \Fz \geq F_{\mathrm{ref}} \label{eq:phi_l}
\end{align}

\begin{theorem}[Gesamtlebensdauer-Maximierung durch kombinierte Systeme]
\label{thm:combined}
Unter den Annahmen \eqref{eq:phi_p}--\eqref{eq:phi_l} gilt:
\begin{equation}
  \Ltire^{\mathrm{combined}} = \frac{h_0 - h_{\min}}{\dot{h}_{\mathrm{basis}}} \cdot
    \frac{1}{\Phi_P^* \cdot \Phi_S^* \cdot \Phi_L^*}
  \label{eq:life_combined}
\end{equation}
Jede Reduzierung der Faktoren $\Phi_P^*, \Phi_S^*, \Phi_L^*$ auf ihren
Minimalwert~1 durch die entsprechenden Steuergeräte verlängert $\Ltire^{\mathrm{combined}}$
monoton und strebt gegen das theoretische Maximum $\Ltire^{\max}$.
\end{theorem}

\begin{proof}
Da alle drei Faktoren $\geq 1$ und paarweise unabhängig sind:
\begin{equation}
  \frac{1}{\Phi_P \Phi_S \Phi_L} \leq 1 \cdot \frac{1}{\min \Phi_P \cdot \min \Phi_S \cdot \min \Phi_L}
    = \frac{1}{1 \cdot 1 \cdot 1} = 1
\end{equation}
Das Optimum $\Phi_i^* = 1$ (für alle $i$) ist gleichzeitig erreichbar:
CRDS setzt $p_i = \pref$ ($\Phi_P = 1$), SCS setzt $\alpha = 0$ ($\Phi_S = 1$),
LCS setzt $\Fz = F_{\mathrm{ref}} = F_{\mathrm{ges}}/4$ ($\Phi_L = 1$).
Damit ist $\Ltire^{\mathrm{combined}} = (h_0 - h_{\min}) / \dot{h}_{\mathrm{basis}}$
maximal.
\end{proof}

\subsection{Quantitative Abschätzung}

Typische Verschleißfaktor-Verbesserungen durch die drei Systeme:

\begin{table}[H]
\centering
\caption{Typische Verschleißfaktoren und erzielte Verbesserungen}
\label{tab:improvement}
\begin{tabular}{lccc}
\toprule
\textbf{System} & \textbf{Ohne Regelung} & \textbf{Mit Regelung} & \textbf{Verbesserung}\\
\midrule
CRDS (Druckregelung)     & $\Phi_P \approx 1{,}25$ & $\Phi_P = 1{,}0$ & $+25\,\%$ \\
SCS (Lenkoptimierung)    & $\Phi_S \approx 1{,}20$ & $\Phi_S = 1{,}0$ & $+20\,\%$ \\
LCS (Lastausgleich)      & $\Phi_L \approx 1{,}35$ & $\Phi_L = 1{,}0$ & $+35\,\%$ \\
\midrule
\textbf{Gesamt (kombiniert)} & $\Phi \approx 2{,}03$ & $\Phi = 1{,}0$ & $\mathbf{+103\,\%}$\\
\bottomrule
\end{tabular}
\end{table}

Das heißt: Im kombinierten Betrieb kann die Reifenlebensdauer theoretisch
mehr als verdoppelt werden. Diese Abschätzung basiert auf den in der Literatur
dokumentierten Einzeleffekten~\cite{Sandberg2011, Niskanen2014, Gent2006}.

\subsection{Systemintegration und CAN-Bus-Kommunikation}

Die drei ECUs kommunizieren über den CAN-Bus des Fahrzeugs (CAN~2.0B,
ISO~11898 \cite{Bosch1991}):

\begin{figure}[H]
\centering
\begin{tabular}{|c|c|c|}
\hline
\textbf{CRDS-ECU} & $\longleftrightarrow$ & \textbf{LCS-ECU}\\
(Druckregelung) & CAN-Bus & (Lastausgleich)\\
\hline
\multicolumn{3}{|c|}{}\\
\multicolumn{3}{|c|}{$\updownarrow$ CAN-Bus}\\
\multicolumn{3}{|c|}{}\\
\hline
\multicolumn{3}{|c|}{\textbf{SCS-ECU (Lenkoptimierung)}}\\
\hline
\end{tabular}
\caption{Vereinfachtes Schema der ECU-Kommunikation über CAN-Bus}
\label{fig:can_schema}
\end{figure}

Relevante CAN-Nachrichten:
\begin{itemize}
  \item CRDS $\to$ SCS/LCS: aktueller Reifendruck, Druckstatus, Temperatur.
  \item LCS $\to$ CRDS: Radlasten (für druckadaptive Korrektur nach Gl.~\eqref{eq:pref_adapt}).
  \item LCS $\to$ SCS: Schwerpunktlage (für Lenkparametrierung).
  \item SCS $\to$ LCS: Querkräfte (für dynamische Lastkorrektur).
\end{itemize}

% ===========================================================
\section{Schlussfolgerungen und Ausblick}
\label{sec:schluss}

\subsection{Zusammenfassung der Ergebnisse}

Diese Arbeit hat formal-mathematisch bewiesen, dass:

\begin{enumerate}
  \item \textbf{CRDS (Theorem~\ref{thm:pressure})}: Ein zentral geregelter
        Reifeninnendruck minimiert die druckbedingte Verschleißrate.
        Die Stabilitätsanalyse (Theorem~\ref{thm:crds_stability}) sichert
        die Realisierbarkeit.
  \item \textbf{SCS (Theorem~\ref{thm:scs_unique}, Theorem~\ref{thm:mpc_wear})}:
        Eine modellprädiktive Lenkungsregelung minimiert die Schräglaufwinkel
        und damit die laterale Verschleißrate.
  \item \textbf{LCS (Theorem~\ref{thm:load_even}, Proposition~\ref{prop:jensen},
        Theorem~\ref{thm:active_suspension})}: Eine gleichmäßige Radlastverteilung,
        unterstützt durch Fahrerhinweise und aktives Fahrwerk, maximiert die
        Lebensdauer des Reifensatzes.
  \item \textbf{Kombination (Theorem~\ref{thm:combined})}: Alle drei Systeme
        wirken multiplikativ und können im kombinierten Betrieb die Reifenlebensdauer
        um über 100\,\% steigern.
\end{enumerate}

\subsection{Praktische Umsetzbarkeit}

Die vorgeschlagenen Systeme sind mit heutiger Fahrzeugtechnik (CAN-Bus,
OBD-II, aktives Fahrwerk, elektrische Servolenkung) vollständig realisierbar.
TPMS (Tyre Pressure Monitoring System) ist in der EU seit 2014 für alle
Neuwagen vorgeschrieben~\cite{EU_TPMS2009}; das CRDS ist eine Erweiterung
hierzu. Modellprädiktive Regler sind in aktuellen Fahrdynamiksystemen
bereits im Einsatz~\cite{Falcone2007}.

\subsection{Ausblick}

Zukünftige Arbeiten könnten:
\begin{itemize}
  \item Digitale Zwillinge des Reifens für eine präzisere Prognose der
        Restlebensdauer (Remaining Useful Life, RUL) einsetzen,
  \item Maschinelles Lernen zur individualisierten Verschleißmodellierung nutzen,
  \item Die Systeme mit V2X-Kommunikation (Fahrzeug-zu-Infrastruktur) koppeln,
        um Fahrbahnzustandsdaten für die Verschleißprognose zu nutzen,
  \item Flotten-übergreifende Reifendatenbanken für statistische Validierung aufbauen.
\end{itemize}

% ===========================================================
\section*{Symbolverzeichnis}
\addcontentsline{toc}{section}{Symbolverzeichnis}

\begin{longtable}{lll}
\toprule
\textbf{Symbol} & \textbf{Einheit} & \textbf{Bedeutung}\\
\midrule
$A_{\mathrm{latsch}}$ & \si{\metre\squared} & Reifenlatschfläche\\
$C_\alpha$ & \si{\newton\per\radian} & Schräglaufsteifigkeit\\
$c_h$ & \si{1\per\metre\per\pascal} & Profiltiefenverschleißkonstante\\
$c_y$ & -- & Seitenverschleißkoeffizient\\
$D$ & \si{\newton} & Maximale Seitenkraft (Pacejka)\\
$E_s$ & \si{\joule} & Schlupfenergie\\
$\Fz$ & \si{\newton} & Radlast\\
$F_{\mathrm{ges}}$ & \si{\newton} & Gesamtfahrzeuglast\\
$f_{\mathrm{dist}}$ & -- & Druckabhängige Kontaktverteilung\\
$h$ & \si{\metre} & Reifenprofiltiefe\\
$h_0$ & \si{\metre} & Anfangsprofiltiefe\\
$h_{\min}$ & \si{\metre} & Mindestprofiltiefe (gesetzlich)\\
$I_z$ & \si{\kilogram\metre\squared} & Massenträgheitsmoment Hochachse\\
$k_E$ & \si{\metre\squared\per\newton} & Energetischer Verschleißkoeffizient\\
$k_w$ & -- & Archard-Verschleißkonstante\\
$l_v, l_h$ & \si{\metre} & Achsabstand vorn/hinten zum Schwerpunkt\\
$L_{\mathrm{tire}}$ & \si{\second} & Reifenlebensdauer\\
$m$ & \si{\kilogram} & Fahrzeugmasse\\
$n$ & -- & Lastexponent im Verschleißmodell\\
$p_i$ & \si{\pascal} & Reifeninnendruck\\
$p_{\mathrm{ref}}$ & \si{\pascal} & Optimaler/Referenz-Reifendruck\\
$r_d$ & \si{\metre} & Dynamischer Reifenrollradius\\
$T_r$ & \si{\celsius} & Reifentemperatur\\
$v$ & \si{\metre\per\second} & Fahrzeuggeschwindigkeit\\
$v_s$ & \si{\metre\per\second} & Schlupfgeschwindigkeit\\
$V_r$ & \si{\metre\cubed} & Reifenvolumen (Luftkammer)\\
$\alpha$ & \si{\radian} & Schräglaufwinkel\\
$\beta$ & \si{\radian} & Schwimmwinkel\\
$\delta$ & \si{\radian} & Lenkwinkel\\
$\dot\psi$ & \si{\radian\per\second} & Gierrate\\
$\kappa$ & -- & Längsschlupf\\
$\sigma_0$ & \si{\pascal} & Mittlerer Kontaktdruck\\
$\phi(\cdot)$ & -- & Konvexe Verschleißfunktion\\
$\Phi_P, \Phi_S, \Phi_L$ & -- & Verschleißfaktoren (Druck, Schlupf, Last)\\
\bottomrule
\end{longtable}

\begin{thebibliography}{99}
\addcontentsline{toc}{section}{Quellenverzeichnis}

\bibitem[Anderson1990]{Anderson1990}
Anderson, B.\,D.\,O.; Moore, J.\,B.: \textit{Optimal Control: Linear Quadratic Methods}.
Prentice Hall, Englewood Cliffs, NJ, 1990.
ISBN~0-13-638651-2.
Inhalt: Standardwerk zur optimalen Regelungstheorie, LQR und LQG-Synthese.
Grundlage für Theorem~\ref{thm:active_suspension}.

\bibitem[Archard1953]{Archard1953}
Archard, J.\,F.: \textit{Contact and Rubbing of Flat Surfaces}.
In: Journal of Applied Physics, Vol.~24, Nr.~8, S.~981--988, 1953.
DOI:~10.1063/1.1721448.
Inhalt: Grundlegende Formulierung des Archard-Verschleißgesetzes (Gl.~\ref{eq:archard}).

\bibitem[Astrom1997]{Astrom1997}
Åström, K.\,J.; Wittenmark, B.: \textit{Computer Controlled Systems: Theory and Design}.
3. Aufl., Prentice Hall, 1997.
ISBN~0-13-314899-8.
Inhalt: Digitale Regelungstechnik, Nyquist- und Hurwitz-Kriterien, Stabilitätsnachweis
für diskrete Regler (Basis für Theorem~\ref{thm:crds_stability}).

\bibitem[Bosch1991]{Bosch1991}
Robert Bosch GmbH: \textit{CAN Specification, Version 2.0}.
Bosch, Stuttgart, 1991.
Inhalt: Spezifikation des CAN-Protokolls (ISO~11898), Grundlage für das
ECU-Kommunikationskonzept in Abschnitt~\ref{sec:gesamtmodell}.

\bibitem[ETRMA2023]{ETRMA2023}
European Tyre and Rubber Manufacturers' Association (ETRMA):
\textit{Annual Statistics -- European Tyre Industry Key Figures 2023}.
ETRMA, Brüssel, 2023. URL: \url{https://www.etrma.org}.
Inhalt: Statistiken zum europäischen Reifenmarkt, Verbrauchszahlen.

\bibitem[EU\_TPMS2009]{EU_TPMS2009}
Europäisches Parlament; Rat der EU:
\textit{Verordnung (EG) Nr.~661/2009 des Europäischen Parlaments und des Rates}
über die Typgenehmigung von Kraftfahrzeugen, Anhang II (TPMS-Anforderungen).
Amtsblatt der Europäischen Union, 2009.
Inhalt: Vorschrift zur Pflicht-Ausstattung von Neuwagen mit Reifendruckkontrollsystemen (TPMS).

\bibitem[Falcone2007]{Falcone2007}
Falcone, P.; Borrelli, F.; Asgari, J.; Tseng, H.\,E.; Hrovat, D.:
\textit{Predictive Active Steering Control for Autonomous Vehicle Systems}.
In: IEEE Transactions on Control Systems Technology, Vol.~15, Nr.~3,
S.~566--580, 2007. DOI:~10.1109/TCST.2007.894653.
Inhalt: Modellprädiktive Lenkungsregelung (MPC) für Fahrzeuge, Basis für
Abschnitt~\ref{sec:scs}.

\bibitem[Gent2006]{Gent2006}
Gent, A.\,N.; Walter, J.\,D. (Hrsg.): \textit{The Pneumatic Tire}.
NHTSA, U.S. Department of Transportation, Washington D.C., 2006.
DOT-HS-810-561.
Inhalt: Umfassendes Standardwerk zur Reifentechnik: Latschgeometrie, Kontaktmechanik,
Verschleiß, Reifendruckeinfluss (Grundlage für Abschnitt~\ref{sec:crds}).

\bibitem[Grosch1963]{Grosch1963}
Grosch, K.\,A.: \textit{The Relation Between the Friction and Visco-Elastic Properties
of Rubber}.
In: Proceedings of the Royal Society of London A, Vol.~274, S.~21--39, 1963.
DOI:~10.1098/rspa.1963.0090.
Inhalt: Visko-elastisches Reibungsmodell für Gummi, Grundlage für energetisches
Verschleißmodell (Gl.~\ref{eq:energetisch}).

\bibitem[Jensen1906]{Jensen1906}
Jensen, J.\,L.\,W.\,V.: \textit{Sur les fonctions convexes et les inégalités entre
les valeurs moyennes}.
In: Acta Mathematica, Vol.~30, S.~175--193, 1906.
DOI:~10.1007/BF02418571.
Inhalt: Mathematische Grundlage der Jensenschen Ungleichung (Proposition~\ref{prop:jensen}).

\bibitem[Kalman1960]{Kalman1960}
Kalman, R.\,E.: \textit{A New Approach to Linear Filtering and Prediction Problems}.
In: Transactions of the ASME -- Journal of Basic Engineering, Vol.~82, S.~35--45, 1960.
DOI:~10.1115/1.3662552.
Inhalt: Kalman-Filter-Algorithmus; Grundlage für die Radlastschätzung
in Abschnitt~\ref{sec:lcs}.

\bibitem[Mayne2000]{Mayne2000}
Mayne, D.\,Q.; Rawlings, J.\,B.; Rao, C.\,V.; Scokaert, P.\,O.\,M.:
\textit{Constrained Model Predictive Control: Stability and Optimality}.
In: Automatica, Vol.~36, Nr.~6, S.~789--814, 2000.
DOI:~10.1016/S0005-1098(99)00214-9.
Inhalt: Standardreferenz für MPC-Stabilität und -Optimalität
(Basis für Theorem~\ref{thm:mpc_wear}).

\bibitem[Niskanen2014]{Niskanen2014}
Niskanen, A.\,J.; Tuononen, A.\,J.:
\textit{Three 3-Axis Accelerometers Fixed on the Inner Liner of a Tyre for
Studying Contact Patch Deformations of the Tyre During Heavy Braking}.
In: Proceedings of the Institution of Mechanical Engineers, Part D:
Journal of Automobile Engineering, Vol.~228, Nr.~10, S.~1157--1167, 2014.
DOI:~10.1177/0954407014526073.
Inhalt: Messung von Reifenkontaktverformungen; empirischer Zusammenhang
zwischen Radlast und Verschleiß.

\bibitem[Pacejka2012]{Pacejka2012}
Pacejka, H.\,B.: \textit{Tire and Vehicle Dynamics}. 3. Aufl.,
Butterworth-Heinemann (Elsevier), Oxford, 2012. ISBN~978-0-08-097016-5.
Inhalt: Standardwerk der Reifendynamik, Magic Formula Modell (Gl.~\ref{eq:magic_formula}),
Latschmodell (Gl.~\ref{eq:latsch}).

\bibitem[Rajamani2006]{Rajamani2006}
Rajamani, R.: \textit{Vehicle Dynamics and Control}.
Springer, New York, 2006. ISBN~978-0-387-26396-0.
Inhalt: Fahrdynamikmodelle (Einspurmodell, Gl.~\ref{eq:bicycle1}--\ref{eq:alpha_h}),
Grundlage für Abschnitt~\ref{sec:scs}.

\bibitem[Sandberg2011]{Sandberg2011}
Sandberg, U.; Bergiers, A.; Ejsmont, J.\,A.; Goubert, L.; Karlsson, R.;
Zöller, M.: \textit{Road Surface Influence on Tyre/Road Rolling Resistance}.
EU-Projekt MIRIAM SP1, 2011. URL: \url{http://www.miriam-co2.net}.
Inhalt: Einfluss von Fahrbahnoberfläche und Reifendruck auf Rollwiderstand
und damit auf Reifenverschleiß.

\bibitem[Schallamach1971]{Schallamach1971}
Schallamach, A.: \textit{How Does Rubber Slide?}
In: Wear, Vol.~17, Nr.~4, S.~301--312, 1971.
DOI:~10.1016/0043-1648(71)90033-0.
Inhalt: Physikalisches Modell für Gleitverhalten und Abrieb von Gummi;
Grundlage des energetischen Verschleißmodells (Gl.~\ref{eq:energetisch}).

\bibitem[Williams1955]{Williams1955}
Williams, M.\,L.; Landel, R.\,F.; Ferry, J.\,D.:
\textit{The Temperature Dependence of Relaxation Mechanisms in Amorphous Polymers
and Other Glass-Forming Liquids}.
In: Journal of the American Chemical Society, Vol.~77, Nr.~14, S.~3701--3707, 1955.
DOI:~10.1021/ja01619a008.
Inhalt: WLF-Gleichung für temperaturabhängige Materialeigenschaften von Polymeren/Gummi
(Gl.~\ref{eq:wlf}).

\end{thebibliography}

\end{document}

